\documentclass[11pt,a4paper]{article}
\usepackage[margin=1in]{geometry}
\usepackage{amsmath,amssymb,amsthm}
\usepackage{algorithm}
\usepackage{algpseudocode}
\usepackage{booktabs}
\usepackage{hyperref}

\theoremstyle{definition}
\newtheorem{definition}{\textbf{Definition}}[section]
\newtheorem{theorem}{\textbf{Theorem}}[section]
\newtheorem{lemma}[theorem]{\textbf{Lemma}}
\newtheorem{remark}{\textbf{Remark}}[section]

\title{\textbf{Latent Diffusion Models}}
\author{\textbf{Team Scaibu} \\ \small Email: \texttt{jaydeep.9664920749@gmail.com} \\ \small \textbf{Architecture and Core Engine Team}}
\date{\textbf{October 2026}}

\begin{document}
\maketitle

\section{Definitions and Cognitive Mental Models}

\subsection{Formal Mathematical Definition}
\textbf{Definition (Two-Stage Perceptual Latent Diffusion Architecture):}
Let $\mathbf{x} \in \mathbb{R}^{H \times W \times 3}$ denote an RGB pixel observation. A Latent Diffusion Model (LDM) decouples generative modeling into two distinct sequential stages:
\begin{enumerate}
    \item \textbf{Perceptual Compression Stage}: A deterministic or regularized autoencoder consisting of an encoder $\mathcal{E}$ and decoder $\mathcal{D}$:
    {\boldmath
    \begin{equation}
    \mathbf{z}_0 = \mathcal{E}(\mathbf{x}) \in \mathbb{R}^{h \times w \times d_{\text{lat}}}, \quad \tilde{\mathbf{x}} = \mathcal{D}(\mathbf{z}_0) \approx \mathbf{x}
    \end{equation}
    }
    where $f = H / h = W / w$ is the spatial downsampling factor (typically $f = 8$). The latent representation is standardized via scaling factor $s > 0$ such that $\mathbf{z}_{\text{scaled}} = s \cdot \mathbf{z}_0$ has empirical identity covariance $\operatorname{Cov}(\mathbf{z}_{\text{scaled}}) \approx \mathbf{I}$.
    \item \textbf{Latent Generative Diffusion Stage}: Forward and reverse diffusion are conducted strictly within the compact latent space:
    {\boldmath
    \begin{equation}
    q(\mathbf{z}_t \mid \mathbf{z}_0) = \mathcal{N}\left(\mathbf{z}_t; \sqrt{\bar{\alpha}_t} \mathbf{z}_{\text{scaled}}, (1 - \bar{\alpha}_t)\mathbf{I}\right)
    \end{equation}
    }
    minimizing the conditional latent denoising score objective:
    {\boldmath
    \begin{equation}
    \mathcal{L}_{\text{LDM}}(\theta) = \mathbb{E}_{\mathbf{z}_0 \sim \mathcal{E}(\mathbf{x}), \boldsymbol{\epsilon} \sim \mathcal{N}(\mathbf{0}, \mathbf{I}), t}\left[ \left\| \boldsymbol{\epsilon} - \boldsymbol{\epsilon}_\theta(\mathbf{z}_t, t, \tau_\phi(y)) \right\|_2^2 \right]
    \end{equation}
    }
    where $\tau_\phi(y)$ is a domain-specific cross-attention conditioning projection (e.g., text, bounding boxes, semantic masks).
\end{enumerate}

\subsection{Expanded Non-Technical Definition (Intuition for Anyone)}
\textbf{Intuitive Conceptual Definition:}
\begin{enumerate}
    \item \textbf{The Heavy Pixel Burden}: Standard diffusion models waste 90\% of their brainpower predicting microscopic pixel noise—like the imperceptible texture of individual cloth fibers or grain on a concrete wall. Computing this on a $1024 \times 1024$ image takes enormous supercomputers.
    \item \textbf{The Master Architect's Blueprint (The Latent Space)}: Instead of working on millions of raw pixels, an ultra-sharp compression lens (Encoder $\mathcal{E}$) squashes the image by $8\times$ in height and width. A giant photo becomes a compact $64 \times 64$ digital blueprint that contains only high-level meaning: shapes, colors, objects, lighting, and composition.
    \item \textbf{Diffusing the Blueprint}: The diffusion model generates only this tiny blueprint! Generating a $64 \times 64$ blueprint is $64\times$ faster and requires vastly less memory.
    \item \textbf{The Master Builder (The Decoder)}: Once the blueprint is fully generated and denoised, the Decoder $\mathcal{D}$ unpacks it in a single quick pass, instantly rendering a stunning, high-resolution full-size image.
\end{enumerate}

\subsection{Child-Friendly Mental Model (Physical Metaphor)}
Imagine building a giant LEGO castle:
\begin{itemize}
    \item \textbf{Pixel Diffusion}: Trying to guess and position every single tiny plastic LEGO speck one by one out of millions of dust particles scattered across your bedroom floor.
    \item \textbf{Latent Diffusion}: Instead of loose dust, you assemble 10 pre-built big LEGO blocks (towers, walls, drawbridges) on a small tabletop. Once your mini-castle is arranged, you push a magic button that instantly expands it into a life-sized real castle!
\end{itemize}

\section{Core Mathematical Equations and Definitions}

\subsection{Spatial Latent Encoding and Variance Standardization}
{\boldmath
\begin{equation}
\mathbf{z}_{\text{scaled}} = s \cdot \mathcal{E}(\mathbf{x}), \quad s = \frac{1}{\sqrt{\frac{1}{N_{\text{lat}} d_{\text{lat}}} \sum_{i, j} (\mathcal{E}(\mathbf{x})[i, j])^2}}
\end{equation}
}
\begin{itemize}
    \item \textbf{Formal Definition}: The affine transformation scaling the raw autoencoder latent distribution to unit variance.
    \item \textbf{What It Means}: Neural networks trained with standard diffusion variance schedules assume the input latent has unit standard deviation. Scaling factor $s$ (e.g. $0.18215$ for Stable Diffusion) aligns the latent code with this prior.
    \item \textbf{Why It Is Important}: Prevents catastrophic gradient explosion or under-diffusion caused by mismatched latent dynamic range.
\end{itemize}

\subsection{Forward Latent Perturbation Equation}
{\boldmath
\begin{equation}
\mathbf{z}_t = \sqrt{\bar{\alpha}_t} \mathbf{z}_{\text{scaled}} + \sqrt{1 - \bar{\alpha}_t} \boldsymbol{\epsilon}, \quad \boldsymbol{\epsilon} \sim \mathcal{N}(\mathbf{0}, \mathbf{I})
\end{equation}
}
\begin{itemize}
    \item \textbf{Formal Definition}: The forward closed-form marginal perturbation operator evaluated in latent coordinates $\mathbb{R}^{h \times w \times d_{\text{lat}}}$.
    \item \textbf{What It Means}: Adds Gaussian noise directly to semantic feature maps instead of raw RGB colors.
    \item \textbf{Why It Is Important}: The diffusion process operates on semantic structure rather than imperceptible high-frequency pixel noise.
\end{itemize}

\subsection{Multimodal Cross-Attention Conditioning Mechanism}
{\boldmath
\begin{equation}
\operatorname{Attention}(\mathbf{Q}, \mathbf{K}, \mathbf{V}) = \operatorname{softmax}\left( \frac{\mathbf{Q} \mathbf{K}^T}{\sqrt{d_k}} \right) \mathbf{V}, \quad \mathbf{Q} = \mathbf{W}_Q^{(i)} \phi_i(\mathbf{z}_t), \quad \mathbf{K} = \mathbf{W}_K^{(i)} \tau_\phi(y), \quad \mathbf{V} = \mathbf{W}_V^{(i)} \tau_\phi(y)
\end{equation}
}
\begin{itemize}
    \item \textbf{Formal Definition}: The cross-attention operator injecting external context $y$ into intermediate U-Net/DiT latent layers $\phi_i(\mathbf{z}_t)$.
    \item \textbf{What It Means}: Spatial latent tokens query linguistic tokens from a pre-trained text encoder (e.g. CLIP or T5) to align image semantics with textual descriptions.
    \item \textbf{Why It Is Important}: Decouples generative visual modeling from text representation learning, allowing arbitrary conditioning modalities (text, depth, pose, audio).
\end{itemize}

\subsection{Spatial Memory Reduction Ratio}
{\boldmath
\begin{equation}
\mathcal{R}_{\text{mem}} = \frac{H \cdot W}{h \cdot w} = f^2
\end{equation}
}
\begin{itemize}
    \item \textbf{Formal Definition}: The quadratic compression factor governing the reduction in spatial grid resolution.
    \item \textbf{What It Means}: For downsampling factor $f = 8$, the spatial resolution drops from $512 \times 512 = 262,144$ pixels to $64 \times 64 = 4,096$ latent tokens, a $64\times$ reduction in spatial memory.
    \item \textbf{Why It Is Important}: Transforms diffusion from an academic supercomputing curiosity into an inference engine that runs on consumer GPUs.
\end{itemize}

\section{Rigorous Mathematical Analysis Checklist}

\subsection{Notation and Symbol Semantics}
\begin{itemize}
    \item $\mathbf{x} \in \mathbb{R}^{H \times W \times 3}$: Clean spatial pixel image.
    \item $\mathbf{z}_0 \in \mathbb{R}^{h \times w \times d_{\text{lat}}}$: Unscaled latent tensor produced by encoder $\mathcal{E}$.
    \item $s \in \mathbb{R}^+$: Latent variance standardization multiplier.
    \item $\bar{\alpha}_t \in (0, 1)$: Cumulative variance schedule at timestep $t$.
    \item $f = H / h$: Spatial downsampling factor ($f \in \{4, 8, 16\}$).
    \item $\mathbf{z}_t \in \mathbb{R}^{h \times w \times d_{\text{lat}}}$: Diffused latent state at timestep $t$.
\end{itemize}

\subsection{Epistemic Status Table}
\begin{table}[h]
\centering
\small
\begin{tabular}{@{}llll@{}}
\toprule
\textbf{Quantity} & \textbf{Symbol} & \textbf{Epistemic Status} & \textbf{Operational Role} \\
\midrule
Perceptual Encoder & $\mathcal{E}$ & Pre-trained Deterministic Map & Compresses high-frequency pixel redundancy \\
Scaling Factor & $s$ & Empirically Calibrated Constant & Normalizes latent variance to unity \\
Diffused Latent & $\mathbf{z}_t$ & Stochastic Intermediate Tensor & Carries noisy semantic representation \\
Perceptual Decoder & $\mathcal{D}$ & Pre-trained Deterministic Map & Expands denoised latent to pixel space \\
\bottomrule
\end{tabular}
\end{table}

\subsection{Computational Dependency Graph}
{\centering
$\mathbf{x} \longrightarrow \mathcal{E}(\mathbf{x}) \longrightarrow \mathbf{z}_0 \longrightarrow s \cdot \mathbf{z}_0 \longrightarrow \{\bar{\alpha}_t, \boldsymbol{\epsilon}\} \longrightarrow \mathbf{z}_t \longrightarrow \text{Latent Diffusion} \longrightarrow \hat{\mathbf{z}}_0 \longrightarrow \mathcal{D}(\hat{\mathbf{z}}_0 / s) \longrightarrow \tilde{\mathbf{x}}$
\par}

\subsection{Dimension and Coordinate Consistency}
Spatial pixel tensors have dimension $H \times W \times 3$; latent tensors have dimension $h \times w \times d_{\text{lat}}$ with $h = H/f, w = W/f$. Cross-attention queries have sequence length $N_{\text{lat}} = h \cdot w$, keys and values have sequence length $N_{\text{text}}$.

\subsection{Asymptotic Regimes}
\begin{itemize}
    \item \textbf{As $f \to 1$}: Recovers pixel-space diffusion; prohibitive memory complexity $\mathcal{O}((HW)^2)$ in attention.
    \item \textbf{As $f \to \infty$}: Severe rate-distortion loss; autoencoder bottleneck drops semantic content, causing blurred decodings.
    \item \textbf{Optimal regime $f \in [4, 8]$}: Balances perceptual compression with maximal computational efficiency.
\end{itemize}

\subsection{Boundary Constraints and Invariants}
\begin{itemize}
    \item \textbf{Reconstruction Invariance}: In the absence of diffusion noise, $\mathcal{D}(\mathcal{E}(\mathbf{x})) \approx \mathbf{x}$ up to perceptual loss bounds $\mathcal{L}_{\text{LPIPS}} < \epsilon$.
\end{itemize}

\subsection{Structural Isomorphisms}
Latent Diffusion is structurally isomorphic to Hierarchical VAEs where the prior distribution over the continuous bottleneck is modeled by a score-based diffusion model.

\section{Theoretical Theorems and Formal Proofs}

\begin{theorem}[Computational Complexity Scaling in Latent Space]
Let a self-attention mechanism operate on an image of spatial dimension $H \times W$. Compression by factor $f = H/h = W/w$ reduces the computational complexity and activation memory of the self-attention operator by a factor of $f^4$:
{\boldmath
\begin{equation}
\frac{\mathcal{C}_{\text{pixel}}}{\mathcal{C}_{\text{latent}}} = f^4
\end{equation}
}
\end{theorem}

\begin{proof}
Consider a standard multi-head self-attention layer. For an input feature map of spatial dimensions $H \times W$, the token sequence length is $N_{\text{pixel}} = H \cdot W$.
The computation of the attention score matrix $\mathbf{A} = \operatorname{softmax}\left(\frac{\mathbf{Q}\mathbf{K}^T}{\sqrt{d}}\right)$ requires computing an inner product between all pairs of tokens:
{\boldmath
\begin{equation}
\mathcal{C}_{\text{pixel}} = 2 \cdot (N_{\text{pixel}})^2 \cdot d = 2 \cdot (H \cdot W)^2 \cdot d
\end{equation}
}
Now consider the compressed latent representation of spatial dimensions $h \times w$ where $h = H/f$ and $w = W/f$. The sequence length is:
{\boldmath
\begin{equation}
N_{\text{latent}} = h \cdot w = \left(\frac{H}{f}\right) \cdot \left(\frac{W}{f}\right) = \frac{H \cdot W}{f^2}
\end{equation}
}
The computational complexity in latent space is:
{\boldmath
\begin{equation}
\mathcal{C}_{\text{latent}} = 2 \cdot (N_{\text{latent}})^2 \cdot d = 2 \cdot \left(\frac{H \cdot W}{f^2}\right)^2 \cdot d = \frac{2 \cdot (H \cdot W)^2 \cdot d}{f^4}
\end{equation}
}
Taking the ratio of the two computational costs:
{\boldmath
\begin{equation}
\frac{\mathcal{C}_{\text{pixel}}}{\mathcal{C}_{\text{latent}}} = \frac{2 (H \cdot W)^2 d}{\frac{2 (H \cdot W)^2 d}{f^4}} = f^4
\end{equation}
}
For a standard downsampling factor of $f = 8$, the attention complexity reduction is $8^4 = 4096\times$, proving the dramatic scaling advantage of latent diffusion.
\end{proof}

\begin{theorem}[Preservation of Diffusion Signal-to-Noise Ratio under Variance Scaling]
If raw encoder latents satisfy $\operatorname{Cov}(\mathbf{z}_0) = \sigma_z^2 \mathbf{I}$, setting scaling factor $s = 1/\sigma_z$ preserves the exact signal-to-noise ratio (SNR) schedule $\operatorname{SNR}(t) = \frac{\bar{\alpha}_t}{1 - \bar{\alpha}_t}$ assumed by standard diffusion priors.
\end{theorem}

\begin{proof}
Let $\mathbf{z}_{\text{scaled}} = s \mathbf{z}_0$. The covariance of the scaled latent is:
{\boldmath
\begin{equation}
\operatorname{Cov}(\mathbf{z}_{\text{scaled}}) = s^2 \operatorname{Cov}(\mathbf{z}_0) = s^2 \sigma_z^2 \mathbf{I}
\end{equation}
}
Setting $s = \frac{1}{\sigma_z}$ yields $\operatorname{Cov}(\mathbf{z}_{\text{scaled}}) = \mathbf{I}$.
Now consider the noisy latent state:
{\boldmath
\begin{equation}
\mathbf{z}_t = \sqrt{\bar{\alpha}_t}\mathbf{z}_{\text{scaled}} + \sqrt{1 - \bar{\alpha}_t}\boldsymbol{\epsilon}, \quad \boldsymbol{\epsilon} \sim \mathcal{N}(\mathbf{0}, \mathbf{I})
\end{equation}
}
The signal variance is $\operatorname{Var}(\sqrt{\bar{\alpha}_t}\mathbf{z}_{\text{scaled}}) = \bar{\alpha}_t \mathbf{I}$.
The noise variance is $\operatorname{Var}(\sqrt{1 - \bar{\alpha}_t}\boldsymbol{\epsilon}) = (1 - \bar{\alpha}_t)\mathbf{I}$.
The resulting signal-to-noise ratio is:
{\boldmath
\begin{equation}
\operatorname{SNR}(t) = \frac{\operatorname{Tr}(\operatorname{Var}(\text{Signal}))}{\operatorname{Tr}(\operatorname{Var}(\text{Noise}))} = \frac{\bar{\alpha}_t d_{\text{lat}}}{(1 - \bar{\alpha}_t) d_{\text{lat}}} = \frac{\bar{\alpha}_t}{1 - \bar{\alpha}_t}
\end{equation}
}
Without variance scaling $s$, the effective SNR would be $\frac{s^2 \bar{\alpha}_t}{1 - \bar{\alpha}_t}$, causing early time-steps to suffer from extreme signal dominance or noise starvation.
\end{proof}

\section{Foundational Architectural Questions and Epistemic Meta-Questions}

\subsection{Architectural Question 1: Perceptual vs Generative Rate-Distortion}
\textbf{Question:} Why separate image synthesis into perceptual compression and generative diffusion rather than training an end-to-end pixel model?\\
\textbf{Answer:} Generative learning suffers when high-capacity neural networks allocate modeling capacity to high-frequency perceptual details imperceptible to humans. Separating compression allows GAN/LPIPS autoencoders to handle imperceptible high frequencies deterministically, leaving diffusion to focus solely on high-entropy semantic content.

\textbf{Meta-Question:} Does two-stage training introduce compounding approximation errors?\\
\textbf{Answer:} Yes; any spatial detail lost by the autoencoder can never be recovered by the diffusion model. However, empirical LPIPS losses with patch discriminators preserve sub-pixel sharpness up to the human perceptual threshold.

\subsection{Architectural Question 2: Spatial Artifacts from Latent Decoding}
\textbf{Question:} Why do decoded images occasionally display tiling or seam artifacts when using spatial latent diffusion?\\
\textbf{Answer:} When decoding high-resolution latents via tiled VAE decoders, boundary discontinuities in overlap blending create visible seams. Utilizing receptive-field-aware border padding and gradient-weighted Gaussian blending across overlapping tiles eliminates seams.

\textbf{Meta-Question:} Can latent diffusion be extended seamlessly to 3D video or audio?\\
\textbf{Answer:} Yes. Temporal or spectral autoencoders compress raw video $(T \times H \times W \times 3)$ or audio waveforms $(T \times 1)$ into compact spatio-temporal or mel-spectrogram latents, where 3D diffusion operates with identical computational savings.

\section{Algorithmic Implementation Specification}

\begin{algorithm}[H]
\caption{Latent Diffusion Spatial Scaling and Marginal Forward Injection}
\begin{algorithmic}[1]
\Require Latent grid $\mathbf{z}_0 \in \mathbb{R}^{N \times D}$, noise grid $\boldsymbol{\epsilon} \in \mathbb{R}^{N \times D}$, cumulative variance $\bar{\alpha}_t \in (0, 1)$, compression factor $f \ge 1$, scaling factor $s > 0$
\Ensure Diffused latent $\mathbf{z}_t \in \mathbb{R}^{N \times D}$, scaled clean latent $\mathbf{z}_{\text{scaled}} \in \mathbb{R}^{N \times D}$, memory reduction ratio $\mathcal{R}_{\text{mem}}$, Frobenius energy $\mathcal{E}$
\State $c_{\text{sig}} \gets \sqrt{\bar{\alpha}_t}$
\State $c_{\text{noise}} \gets \sqrt{1 - \bar{\alpha}_t}$
\State $\mathcal{E} \gets 0.0$
\State Initialize empty arrays $\mathbf{z}_{\text{scaled}}, \mathbf{z}_t$ of size $N \times D$
\For{$i = 0$ \textbf{to} $N - 1$}
    \For{$j = 0$ \textbf{to} $D - 1$}
        \State $v_{\text{scaled}} \gets \mathbf{z}_0[i][j] \cdot s$
        \State $\mathbf{z}_{\text{scaled}}[i][j] \gets v_{\text{scaled}}$
        \State $v_t \gets c_{\text{sig}} \cdot v_{\text{scaled}} + c_{\text{noise}} \cdot \boldsymbol{\epsilon}[i][j]$
        \State $\mathbf{z}_t[i][j] \gets v_t$
        \State $\mathcal{E} \gets \mathcal{E} + v_t \cdot v_t$
    \EndFor
\EndFor
\State $\mathcal{R}_{\text{mem}} \gets f \cdot f$
\State \Return $(\mathbf{z}_t, \mathbf{z}_{\text{scaled}}, \mathcal{R}_{\text{mem}}, \mathcal{E})$
\end{algorithmic}
\end{algorithm}

\section{Big-$\Theta$ Complexity Analysis}

\begin{itemize}
    \item \textbf{Latent Injection Time Complexity}: $\Theta(N_{\text{lat}} \cdot d_{\text{lat}})$ scalar operations where $N_{\text{lat}} = h \cdot w$.
    \item \textbf{Decoded Image Synthesis Time Complexity}: $\Theta(N_{\text{lat}} \cdot d_{\text{lat}} + \mathcal{C}_{\text{VAE\_Dec}})$, where the decoder runs in a single forward pass.
    \item \textbf{Space Complexity}: $\Theta(N_{\text{lat}} \cdot d_{\text{lat}})$ memory, which is $f^2\times$ smaller than raw pixel buffers.
    \item \textbf{Work \& Depth}: Work is $\mathcal{W} = \Theta(N_{\text{lat}} \cdot d_{\text{lat}})$; parallel depth is $\mathcal{D} = \Theta(\log(N_{\text{lat}} \cdot d_{\text{lat}}))$.
\end{itemize}

\section{Single Canonical Repository Reference}

The complete deterministic scalar implementation, verification suite, and unit test benchmarks are published at:
\begin{center}
\url{https://github.com/Chief-Strategist-J/llm-obs-infra}
\end{center}

\end{document}
